\nonstopmode{}
\documentclass[letterpaper]{book}
\usepackage[times,inconsolata,hyper]{Rd}
\usepackage{makeidx}
\usepackage[utf8]{inputenc} % @SET ENCODING@
% \usepackage{graphicx} % @USE GRAPHICX@
\makeindex{}
\begin{document}
\chapter*{}
\begin{center}
{\textbf{\huge Package `AlphaSDM'}}
\par\bigskip{\large \today}
\end{center}
\inputencoding{utf8}
\ifthenelse{\boolean{Rd@use@hyper}}{\hypersetup{pdftitle = {AlphaSDM: Species Distribution Modeling using Alpha Earth Embeddings}}}{}
\begin{description}
\raggedright{}
\item[Title]\AsIs{Species Distribution Modeling using Alpha Earth Embeddings}
\item[Version]\AsIs{0.1.0}
\item[Author]\AsIs{James Longo [aut, cre]}
\item[Maintainer]\AsIs{James Longo }\email{james.longo.birds@gmail.com}\AsIs{}
\item[Description]\AsIs{Facilitates Species Distribution Modeling (SDM) by leveraging Google's Alpha Earth satellite embeddings. The package provides tools for extracting high-resolution (10m) 64-dimensional dense vectors from Google Earth Engine (GEE), training niche centroid and machine learning models, and generating suitability maps. All computationally intensive tasks are performed server-side on GEE to ensure scalability for large datasets.}
\item[License]\AsIs{MIT + file LICENSE}
\item[Encoding]\AsIs{UTF-8}
\item[RoxygenNote]\AsIs{7.3.3}
\item[Imports]\AsIs{parallel, reticulate, rgee, jsonlite, sf, stars, data.table,
magrittr}
\item[URL]\AsIs{}\url{https://github.com/James-Longo/AlphaSDM}\AsIs{}
\item[BugReports]\AsIs{}\url{https://github.com/James-Longo/AlphaSDM/issues}\AsIs{}
\item[Suggests]\AsIs{testthat (>= 3.0.0), knitr, rmarkdown}
\item[VignetteBuilder]\AsIs{knitr}
\item[Config/testthat/edition]\AsIs{3}
\item[NeedsCompilation]\AsIs{no}
\end{description}
\Rdcontents{\R{} topics documented:}
\inputencoding{utf8}
\HeaderA{.run\_command\_with\_timer}{Shared internal helper for live-updating timer Shared internal helper for live-updating timer}{.run.Rul.command.Rul.with.Rul.timer}
%
\begin{Description}
Shared internal helper for live-updating timer
Shared internal helper for live-updating timer
\end{Description}
%
\begin{Usage}
\begin{verbatim}
.run_command_with_timer(label, active = FALSE)
\end{verbatim}
\end{Usage}
%
\begin{Arguments}
\begin{ldescription}
\item[\code{label}] Text label.

\item[\code{active}] Boolean logic.
\end{ldescription}
\end{Arguments}
\inputencoding{utf8}
\HeaderA{analyze\_embeddings}{Analyze Embeddings}{analyze.Rul.embeddings}
\keyword{internal}{analyze\_embeddings}
%
\begin{Description}
This function calculates the similarity between each location's embedding and the species centroid (mean embedding).
\end{Description}
%
\begin{Usage}
\begin{verbatim}
analyze_embeddings(df, method, gee_project = NULL, cv = FALSE, scale)
\end{verbatim}
\end{Usage}
%
\begin{Arguments}
\begin{ldescription}
\item[\code{df}] Data frame containing columns A00 to A63.

\item[\code{method}] Mapping method: 'centroid' or 'ridge'.
\end{ldescription}
\end{Arguments}
%
\begin{Value}
A list containing the results.
\end{Value}
\inputencoding{utf8}
\HeaderA{assign\_spatial\_folds}{Assign Spatial Folds}{assign.Rul.spatial.Rul.folds}
\keyword{internal}{assign\_spatial\_folds}
%
\begin{Description}
Assign Spatial Folds
\end{Description}
%
\begin{Usage}
\begin{verbatim}
assign_spatial_folds(df, n_folds = 10)
\end{verbatim}
\end{Usage}
%
\begin{Arguments}
\begin{ldescription}
\item[\code{df}] Data frame with longitude, latitude

\item[\code{n\_folds}] Number of folds
\end{ldescription}
\end{Arguments}
\inputencoding{utf8}
\HeaderA{calculate\_cbi}{Calculate Continuous Boyce Index (CBI)}{calculate.Rul.cbi}
%
\begin{Description}
Calculate Continuous Boyce Index (CBI)
\end{Description}
%
\begin{Usage}
\begin{verbatim}
calculate_cbi(pos_scores, all_scores, window_width = 0.1, n_bins = 100)
\end{verbatim}
\end{Usage}
%
\begin{Arguments}
\begin{ldescription}
\item[\code{pos\_scores}] Scores for presence points

\item[\code{all\_scores}] Scores for all points (presence + background)

\item[\code{window\_width}] Width of the moving window (proportion of score range)

\item[\code{n\_bins}] Number of bins for evaluation
\end{ldescription}
\end{Arguments}
\inputencoding{utf8}
\HeaderA{calculate\_classifier\_metrics}{Calculate Classifier Metrics}{calculate.Rul.classifier.Rul.metrics}
%
\begin{Description}
Calculate Classifier Metrics
\end{Description}
%
\begin{Usage}
\begin{verbatim}
calculate_classifier_metrics(scores_pos, scores_neg)
\end{verbatim}
\end{Usage}
%
\begin{Arguments}
\begin{ldescription}
\item[\code{scores\_pos}] Scores for presence points

\item[\code{scores\_neg}] Scores for background/absence points
\end{ldescription}
\end{Arguments}
\inputencoding{utf8}
\HeaderA{ee\_hard\_reset}{Hard Reset Google Earth Engine Credentials}{ee.Rul.hard.Rul.reset}
%
\begin{Description}
This function forcefully removes all local Google Earth Engine credentials
to resolve persistent authentication loops or "expired token" errors.
After running this, your next GEE-related call will prompt for a fresh browser authentication.
\end{Description}
%
\begin{Usage}
\begin{verbatim}
ee_hard_reset()
\end{verbatim}
\end{Usage}
\inputencoding{utf8}
\HeaderA{ensure\_AlphaSDM\_dependencies}{Ensure dependencies are met (Internal)}{ensure.Rul.AlphaSDM.Rul.dependencies}
\keyword{internal}{ensure\_AlphaSDM\_dependencies}
%
\begin{Description}
Ensure dependencies are met (Internal)
\end{Description}
%
\begin{Usage}
\begin{verbatim}
ensure_AlphaSDM_dependencies()
\end{verbatim}
\end{Usage}
\inputencoding{utf8}
\HeaderA{ensure\_gee\_authenticated}{Ensure Google Earth Engine is Authenticated and Python is Setup (Internal)}{ensure.Rul.gee.Rul.authenticated}
\keyword{internal}{ensure\_gee\_authenticated}
%
\begin{Description}
Ensure Google Earth Engine is Authenticated and Python is Setup (Internal)
\end{Description}
%
\begin{Usage}
\begin{verbatim}
ensure_gee_authenticated(project = NULL)
\end{verbatim}
\end{Usage}
\inputencoding{utf8}
\HeaderA{evaluate\_models}{Evaluate SDM Models}{evaluate.Rul.models}
%
\begin{Description}
Evaluate SDM Models
\end{Description}
%
\begin{Usage}
\begin{verbatim}
evaluate_models(
  data,
  predict_coords,
  scale = 10,
  output_dir = getwd(),
  methods = NULL,
  aoi_year = NULL,
  count = NULL,
  n_trees = 100L,
  min_leaf_population = 5L,
  bag_fraction = 0.5,
  shrinkage = 0.005,
  max_nodes = NULL,
  variables_per_split = NULL,
  lambda_ = 0,
  polynomial = 2L,
  gee_project = NULL,
  python_path = NULL,
  options = list()
)
\end{verbatim}
\end{Usage}
%
\begin{Arguments}
\begin{ldescription}
\item[\code{data}] A data frame formatted via `format\_data()`.

\item[\code{predict\_coords}] Data frame of target coordinates.

\item[\code{scale}] Resolution in meters.

\item[\code{output\_dir}] Directory to save results.

\item[\code{methods}] Character vector of method names.

\item[\code{aoi\_year}] Alpha Earth Mosaic year.

\item[\code{count}] Number of background points.

\item[\code{n\_trees}] Number of trees for tree-based methods.

\item[\code{min\_leaf\_population}] Min leaf size for trees.

\item[\code{bag\_fraction}] Bagging fraction.

\item[\code{shrinkage}] Learning rate for GBT.

\item[\code{max\_nodes}] Max nodes per tree.

\item[\code{variables\_per\_split}] Variables per split.

\item[\code{lambda\_}] Penalty for Ridge.

\item[\code{polynomial}] Degree of polynomial terms for Ridge regression.

\item[\code{gee\_project}] Optional. Google Cloud Project ID.

\item[\code{python\_path}] Optional. Path to Python executable.

\item[\code{options}] Named list of advanced hyperparameters.
\end{ldescription}
\end{Arguments}
\inputencoding{utf8}
\HeaderA{extrapolate}{Generate Niche Similarity Map}{extrapolate}
\keyword{internal}{extrapolate}
%
\begin{Description}
This function generates a continuous niche similarity map (dot product) in Google Earth Engine
based on the species centroid. The map is returned at the specified resolution.
\end{Description}
%
\begin{Usage}
\begin{verbatim}
extrapolate(
  df,
  analysis_meta_path,
  output_dir = getwd(),
  scale = NULL,
  aoi_year = NULL,
  view = FALSE,
  gcs_bucket = NULL,
  wait = FALSE,
  zip = FALSE,
  python_path = NULL,
  gee_project = NULL
)
\end{verbatim}
\end{Usage}
%
\begin{Arguments}
\begin{ldescription}
\item[\code{df}] Data frame containing the original survey data (Latitude, Longitude) to determine the AOI.

\item[\code{analysis\_meta\_path}] Path to the JSON metadata file generated by `analyze\_embeddings`.

\item[\code{output\_dir}] Directory where the output map and metadata will be saved. Defaults to current working directory.

\item[\code{scale}] Resolution in meters for the output map. Defaults to 10.

\item[\code{gcs\_bucket}] Optional. GCS bucket name for server-side export. If provided, bypasses local download.

\item[\code{wait}] Logical. If TRUE, waits for the GCS export task(s) to complete and shows progress updates.

\item[\code{zip}] Logical. If TRUE, zips the output rasters into a single archive (local mode only).

\item[\code{coarse\_meta\_path}] Optional. Path to metadata for 1km filtering. If provided, analysis is restricted to 1km pixels > 5th percentile.
\end{ldescription}
\end{Arguments}
%
\begin{Value}
A list containing the paths to the generated map files, HTML, or GCS export details.
\end{Value}
\inputencoding{utf8}
\HeaderA{format\_data}{Format Data for AlphaSDM}{format.Rul.data}
%
\begin{Description}
This function standardizes input data frames for use in the AlphaSDM pipeline.
It renames coordinate/year/presence columns to standard lowercase names,
preserves Alpha Earth embeddings if present, and removes all other columns.
\end{Description}
%
\begin{Usage}
\begin{verbatim}
format_data(data, coords, year, presence = NULL, species = NULL)
\end{verbatim}
\end{Usage}
%
\begin{Arguments}
\begin{ldescription}
\item[\code{data}] A data frame containing your survey data.

\item[\code{coords}] A character vector of length 2 specifying the longitude and latitude columns IN ORDER: c(longitude\_col, latitude\_col). Note: Longitude first, then Latitude!

\item[\code{year}] A character string specifying the year or date column.

\item[\code{presence}] Optional. A character string specifying the presence/absence column (values should be 0 or 1).

\item[\code{species}] Optional. A character string specifying the species name column.
\end{ldescription}
\end{Arguments}
%
\begin{Value}
A standardized data frame ready for `AlphaSDM()` with lowercase column names.
\end{Value}
\inputencoding{utf8}
\HeaderA{GEE\_CLASSIFIER\_METHODS}{GEE Classifier Methods Registry}{GEE.Rul.CLASSIFIER.Rul.METHODS}
\keyword{internal}{GEE\_CLASSIFIER\_METHODS}
%
\begin{Description}
GEE Classifier Methods Registry
\end{Description}
%
\begin{Usage}
\begin{verbatim}
GEE_CLASSIFIER_METHODS
\end{verbatim}
\end{Usage}
%
\begin{Format}
An object of class \code{list} of length 4.
\end{Format}
\inputencoding{utf8}
\HeaderA{GEE\_REDUCER\_METHODS}{GEE Reducer Methods Registry}{GEE.Rul.REDUCER.Rul.METHODS}
\keyword{internal}{GEE\_REDUCER\_METHODS}
%
\begin{Description}
GEE Reducer Methods Registry
\end{Description}
%
\begin{Usage}
\begin{verbatim}
GEE_REDUCER_METHODS
\end{verbatim}
\end{Usage}
%
\begin{Format}
An object of class \code{list} of length 2.
\end{Format}
\inputencoding{utf8}
\HeaderA{generate\_map}{Generate SDM Map}{generate.Rul.map}
%
\begin{Description}
Trains models on provided data and generates spatial prediction maps (rasters) on GEE.
\end{Description}
%
\begin{Usage}
\begin{verbatim}
generate_map(
  data,
  aoi,
  scale = 10,
  output_dir = getwd(),
  methods = NULL,
  ensemble = TRUE,
  aoi_year = NULL,
  count = NULL,
  n_trees = 100L,
  min_leaf_population = 5L,
  bag_fraction = 0.5,
  shrinkage = 0.005,
  max_nodes = NULL,
  variables_per_split = NULL,
  lambda_ = 0,
  polynomial = 2L,
  gee_project = NULL,
  python_path = NULL,
  options = list()
)
\end{verbatim}
\end{Usage}
%
\begin{Arguments}
\begin{ldescription}
\item[\code{data}] A data frame formatted via `format\_data()`. Must have (longitude, latitude, year, present).
If a 'species' column is present with multiple species, the function processes each separately.

\item[\code{aoi}] Either a list with `lat`, `lon`, and `radius` (in meters), or a path to a polygon file.

\item[\code{scale}] Resolution in meters for the final map. Defaults to 10.

\item[\code{output\_dir}] Directory to save results. Defaults to the current working directory.

\item[\code{methods}] Character vector of method names. Supported: "rf", "gbt", "maxent", "svm", "ridge", "centroid".

\item[\code{ensemble}] Combine all methods into an ensemble map. Defaults to TRUE.

\item[\code{aoi\_year}] Alpha Earth Mosaic year for map generation. Defaults to current year.

\item[\code{count}] Number of background points. Defaults to 10x presence points.

\item[\code{n\_trees}] Number of trees for tree-based methods.

\item[\code{min\_leaf\_population}] Min leaf size for trees.

\item[\code{bag\_fraction}] Bagging fraction.

\item[\code{shrinkage}] Learning rate for GBT.

\item[\code{max\_nodes}] Max nodes per tree.

\item[\code{variables\_per\_split}] Variables per split.

\item[\code{lambda\_}] Penalty for Ridge.

\item[\code{polynomial}] Degree of polynomial terms for Ridge regression. 1 = linear, 2 = quadratic (default).

\item[\code{gee\_project}] Optional. Google Cloud Project ID.

\item[\code{python\_path}] Optional. Path to Python executable.

\item[\code{options}] Named list of advanced hyperparameters.
\end{ldescription}
\end{Arguments}
\inputencoding{utf8}
\HeaderA{get\_background\_embeddings\_gee}{Generate Background Embeddings}{get.Rul.background.Rul.embeddings.Rul.gee}
\keyword{internal}{get\_background\_embeddings\_gee}
%
\begin{Description}
Generate Background Embeddings
\end{Description}
%
\begin{Usage}
\begin{verbatim}
get_background_embeddings_gee(data, scale, count)
\end{verbatim}
\end{Usage}
%
\begin{Arguments}
\begin{ldescription}
\item[\code{data}] Training data frame containing longitude, latitude, and year

\item[\code{scale}] Resolution in meters

\item[\code{count}] Total number of points to generate
\end{ldescription}
\end{Arguments}
\inputencoding{utf8}
\HeaderA{get\_embeddings\_at\_fc}{Sample Embeddings at FeatureCollection}{get.Rul.embeddings.Rul.at.Rul.fc}
\keyword{internal}{get\_embeddings\_at\_fc}
%
\begin{Description}
Sample Embeddings at FeatureCollection
\end{Description}
%
\begin{Usage}
\begin{verbatim}
get_embeddings_at_fc(
  fc,
  scale,
  properties = NULL,
  geometries = FALSE,
  years = NULL
)
\end{verbatim}
\end{Usage}
%
\begin{Arguments}
\begin{ldescription}
\item[\code{fc}] FeatureCollection with 'year' property

\item[\code{scale}] Resolution in meters

\item[\code{properties}] Optional properties to retain

\item[\code{geometries}] Boolean, retain geometries?

\item[\code{years}] Optional list of years to filter
\end{ldescription}
\end{Arguments}
\inputencoding{utf8}
\HeaderA{install\_AlphaSDM}{Install AlphaSDM Python Dependencies}{install.Rul.AlphaSDM}
%
\begin{Description}
This function installs the required Python dependencies for AlphaSDM
using reticulate.
\end{Description}
%
\begin{Usage}
\begin{verbatim}
install_AlphaSDM(method = "virtualenv", envname = "r-AlphaSDM")
\end{verbatim}
\end{Usage}
%
\begin{Arguments}
\begin{ldescription}
\item[\code{method}] Installation method. Pass "auto" to let reticulate decide.

\item[\code{envname}] The name, or full path, of the environment in which Python packages are to be installed.
\end{ldescription}
\end{Arguments}
\inputencoding{utf8}
\HeaderA{predict\_all\_models\_gee}{Predict Scores for Multiple Models on GEE}{predict.Rul.all.Rul.models.Rul.gee}
\keyword{internal}{predict\_all\_models\_gee}
%
\begin{Description}
Consolidates scoring to minimize GEE map passes.
\end{Description}
%
\begin{Usage}
\begin{verbatim}
predict_all_models_gee(fc, models_list)
\end{verbatim}
\end{Usage}
%
\begin{Arguments}
\begin{ldescription}
\item[\code{fc}] GEE FeatureCollection with embeddings

\item[\code{models\_list}] List of model results from AlphaSDM
\end{ldescription}
\end{Arguments}
\inputencoding{utf8}
\HeaderA{predict\_at\_coords}{Predict at Specific Coordinates}{predict.Rul.at.Rul.coords}
%
\begin{Description}
Predict at Specific Coordinates
\end{Description}
%
\begin{Usage}
\begin{verbatim}
predict_at_coords(
  df,
  analysis_meta_paths,
  scale = NULL,
  aoi_year = NULL,
  gee_project = NULL,
  options = list()
)
\end{verbatim}
\end{Usage}
%
\begin{Arguments}
\begin{ldescription}
\item[\code{df}] Data frame with columns `longitude`, `latitude`, and `year`.

\item[\code{analysis\_meta\_paths}] Vector of paths to metadata JSON files.

\item[\code{scale}] Resolution in meters.

\item[\code{aoi\_year}] Year for Alpha Earth mosaic.

\item[\code{gee\_project}] GEE Project ID.

\item[\code{options}] Named list of advanced hyperparameters.
\end{ldescription}
\end{Arguments}
%
\begin{Value}
A list with data (df + similarity) and metrics.
\end{Value}
\inputencoding{utf8}
\HeaderA{predict\_gee\_map}{Predict GEE Map}{predict.Rul.gee.Rul.map}
\keyword{internal}{predict\_gee\_map}
%
\begin{Description}
Predict GEE Map
\end{Description}
%
\begin{Usage}
\begin{verbatim}
predict_gee_map(model_res, img)
\end{verbatim}
\end{Usage}
%
\begin{Arguments}
\begin{ldescription}
\item[\code{model\_res}] Result from train\_gee\_model

\item[\code{img}] Alpha Earth mosaic
\end{ldescription}
\end{Arguments}
\inputencoding{utf8}
\HeaderA{predict\_points\_gee}{Predict Scores for Single Model on GEE}{predict.Rul.points.Rul.gee}
\keyword{internal}{predict\_points\_gee}
%
\begin{Description}
Predict Scores for Single Model on GEE
\end{Description}
%
\begin{Usage}
\begin{verbatim}
predict_points_gee(fc, model_res)
\end{verbatim}
\end{Usage}
%
\begin{Arguments}
\begin{ldescription}
\item[\code{fc}] GEE FeatureCollection with embeddings

\item[\code{model\_res}] Model result
\end{ldescription}
\end{Arguments}
\inputencoding{utf8}
\HeaderA{prep\_training\_data\_gee}{Prepare Training Data for GEE Prepare Training Data on GEE (Pure Server-Side)}{prep.Rul.training.Rul.data.Rul.gee}
\keyword{internal}{prep\_training\_data\_gee}
%
\begin{Description}
Prepare Training Data for GEE
Prepare Training Data on GEE (Pure Server-Side)
\end{Description}
%
\begin{Usage}
\begin{verbatim}
prep_training_data_gee(df, class_property = "present", scale = 10)
\end{verbatim}
\end{Usage}
%
\begin{Arguments}
\begin{ldescription}
\item[\code{df}] Data frame with longitude, latitude, year, and class\_property

\item[\code{class\_property}] Column name for presence/absence

\item[\code{scale}] Resolution in meters
\end{ldescription}
\end{Arguments}
\inputencoding{utf8}
\HeaderA{resolve\_python\_path}{Resolve Python Path (Internal)}{resolve.Rul.python.Rul.path}
\keyword{internal}{resolve\_python\_path}
%
\begin{Description}
Resolve Python Path (Internal)
\end{Description}
%
\begin{Usage}
\begin{verbatim}
resolve_python_path(path = NULL)
\end{verbatim}
\end{Usage}
\inputencoding{utf8}
\HeaderA{retry\_curl\_download}{Internal helper for retrying GEE operations with exponential backoff}{retry.Rul.curl.Rul.download}
\keyword{internal}{retry\_curl\_download}
%
\begin{Description}
Internal helper for retrying GEE operations with exponential backoff
\end{Description}
%
\begin{Usage}
\begin{verbatim}
retry_curl_download(expr, max_retries = 5, initial_delay = 1)
\end{verbatim}
\end{Usage}
\inputencoding{utf8}
\HeaderA{timestamp\_message}{Centralized logging with timestamps}{timestamp.Rul.message}
\keyword{internal}{timestamp\_message}
%
\begin{Description}
Centralized logging with timestamps
\end{Description}
%
\begin{Usage}
\begin{verbatim}
timestamp_message(msg)
\end{verbatim}
\end{Usage}
%
\begin{Arguments}
\begin{ldescription}
\item[\code{msg}] Message string
\end{ldescription}
\end{Arguments}
\inputencoding{utf8}
\HeaderA{train\_AlphaSDM\_internal}{Internal Training Pipeline}{train.Rul.AlphaSDM.Rul.internal}
\keyword{internal}{train\_AlphaSDM\_internal}
%
\begin{Description}
Internal Training Pipeline
\end{Description}
%
\begin{Usage}
\begin{verbatim}
train_AlphaSDM_internal(
  data,
  methods,
  aoi_geom = NULL,
  scale = 10,
  aoi_year = NULL,
  count = NULL,
  training_params = list()
)
\end{verbatim}
\end{Usage}
\inputencoding{utf8}
\HeaderA{train\_gee\_model}{Train GEE Model}{train.Rul.gee.Rul.model}
\keyword{internal}{train\_gee\_model}
%
\begin{Description}
Train GEE Model
\end{Description}
%
\begin{Usage}
\begin{verbatim}
train_gee_model(
  sampled_fc,
  method,
  params = list(),
  class_property = "present"
)
\end{verbatim}
\end{Usage}
%
\begin{Arguments}
\begin{ldescription}
\item[\code{sampled\_fc}] Sampled FeatureCollection (from get\_embeddings\_at\_fc)

\item[\code{method}] Method name

\item[\code{params}] List of hyperparameters

\item[\code{class\_property}] Column name for labels
\end{ldescription}
\end{Arguments}
\inputencoding{utf8}
\HeaderA{upload\_points\_to\_gee}{Upload Individual Points to GEE with IDs}{upload.Rul.points.Rul.to.Rul.gee}
\keyword{internal}{upload\_points\_to\_gee}
%
\begin{Description}
Upload Individual Points to GEE with IDs
\end{Description}
%
\begin{Usage}
\begin{verbatim}
upload_points_to_gee(df, chunk_size = 2000, id_offset = 0)
\end{verbatim}
\end{Usage}
%
\begin{Arguments}
\begin{ldescription}
\item[\code{df}] Data frame with longitude, latitude, year

\item[\code{chunk\_size}] Points per upload chunk

\item[\code{id\_offset}] Numeric offset for tmp\_id to support outer batching
\end{ldescription}
\end{Arguments}
%
\begin{Value}
GEE FeatureCollection
\end{Value}
\printindex{}
\end{document}
